\documentclass[11pt]{article}

\usepackage[margin=1in]{geometry}
\usepackage[T1]{fontenc}
\usepackage[utf8]{inputenc}
\usepackage{microtype}
\usepackage{booktabs}
\usepackage{array}
\usepackage{hyperref}
\usepackage[numbers,sort&compress]{natbib}

\title{Workplace Determinants of Metabolic Syndrome in Working Adults: A Systematic Review}
\author{Gideon Ofosu Kyere\\
Kwame Nkrumah University of Science and Technology (KNUST), Kumasi, Ghana\\
ORCID: 0009-0003-9848-8437}
\date{}

\begin{document}
\maketitle

\begin{abstract}
\textbf{Objectives:} To synthesise evidence on occupational stress, diet, physical activity, sedentary behaviour and workplace lifestyle interventions in relation to metabolic syndrome (MetS) among working adults.

\textbf{Methods:} Peer-reviewed studies published from 2015 to 2026 were identified through a workplace-focused PubMed/MEDLINE search supplemented by targeted searches, reference screening and citation chasing. Eligible studies involved employed adults or occupational populations and examined workplace stress, diet, physical activity, sedentary behaviour or workplace lifestyle interventions in relation to MetS or established MetS components. Randomised trials, cohort studies, analytical cross-sectional studies and quasi-experimental interventions were eligible. Design-specific Joanna Briggs Institute tools were used for critical appraisal. Because of substantial clinical and methodological heterogeneity, findings were synthesised narratively.

\textbf{Results:} The current synthesis comprised 34 primary studies. Prospective studies linked high or increasing work stress with subsequent MetS risk, whereas cross-sectional stress findings were less consistent. Higher leisure-time physical activity was generally associated with lower MetS risk, while prolonged sedentary exposure was associated with higher risk. Healthier dietary patterns were associated with more favourable metabolic profiles. Workplace interventions involving exercise, dietary modification, counselling or multicomponent lifestyle support improved waist circumference, blood pressure, glucose regulation, lipid measures or MetS severity in several studies.

\textbf{Conclusions:} MetS in working adults is associated with a combination of occupational and behavioural factors. The strongest evidence supports adverse work stress, physical inactivity and prolonged sedentary behaviour as relevant risk markers, while workplace lifestyle interventions can improve several metabolic outcomes. More consistent longitudinal and randomised evidence is needed across diverse occupational settings.
\end{abstract}

\section*{What is already known on this topic}
Occupational stress, diet, physical inactivity and sedentary behaviour have each been linked to cardiometabolic risk in workers, but the evidence is dispersed across occupational groups and study designs.

\section*{What this study adds}
This review integrates prospective, cross-sectional and workplace-intervention evidence across four domains: occupational stress, diet/nutrition, physical activity/sedentary behaviour and multicomponent workplace programmes. Prospective evidence was strongest for work stress, while intervention studies showed that workplace-based lifestyle modification can improve several MetS components.

\section*{How this study might affect research, practice or policy}
Workplace prevention may benefit from combining psychosocial risk management, healthier food environments, opportunities for physical activity, reduced prolonged sitting and targeted cardiometabolic screening. Future research should prioritise standardised MetS definitions and stronger longitudinal and randomised designs.

\section{Introduction}

Metabolic syndrome (MetS) brings together a cluster of cardiometabolic abnormalities---typically central obesity, elevated blood pressure, impaired glucose regulation, raised triglycerides and reduced high-density lipoprotein cholesterol---that substantially increase the risk of type 2 diabetes and cardiovascular disease.

The workplace is an important setting in which these risks accumulate. Many adults spend a large proportion of waking hours at work, where physical demands, sitting time, shift patterns, psychosocial pressure and access to food can shape everyday behaviour. A sedentary job may reduce energy expenditure; irregular or extended hours may disrupt eating and recovery; and sustained psychosocial stress may influence both behaviour and physiological regulation. At the same time, the workplace offers a practical platform for prevention through screening, food-environment changes, exercise programmes, counselling and digital health support.

The occupational evidence is heterogeneous. Prospective studies have linked high or increasing work stress with later MetS \citep{Garbarino2015,Yamaguchi2018}, while large workforce cohorts have shown differences in MetS risk across occupational groups \citep{vanZon2020Occupational,Runge2021Metabolic}. Leisure-time exercise appears protective in some worker cohorts \citep{Kuwahara2016}, whereas prolonged sedentary exposure has been associated with higher MetS prevalence \citep{Motuma2020}. Intervention studies further suggest that workplace-based exercise, dietary modification and multicomponent lifestyle programmes can improve metabolic outcomes \citep{Haufe2019,Ma2025,Bogale2025}.

What is less clear is how these strands of evidence fit together. Studies differ in occupation, exposure definition, intervention intensity, MetS criteria and methodological design, making a simple pooled estimate inappropriate. This review therefore synthesised the evidence across occupational stress, diet and workplace nutrition, physical activity and sedentary behaviour, and workplace lifestyle interventions among working adults.

\section{Methods}

\subsection{Review design and eligibility}
This review was developed from an earlier undergraduate systematic review, but the journal manuscript used a revised workplace-focused eligibility framework and a fresh literature search. Adults in employment, occupational groups or explicitly defined workplace settings were eligible. Exposures and interventions included occupational stress or job strain, dietary patterns and nutrition, physical activity or exercise, sedentary behaviour, and workplace lifestyle interventions.

Studies had to report MetS as a defined syndrome or report established MetS components in a context explicitly concerned with MetS risk, prevention, severity or management. Eligible designs were randomised controlled trials, prospective or longitudinal cohort studies, analytical cross-sectional studies and quasi-experimental workplace interventions. Protocols, editorials, conference abstracts, grey literature and secondary reviews were excluded. The publication window was 2015--2026.

\subsection{Search and study selection}
PubMed/MEDLINE was the primary bibliographic source. Search concepts combined MetS with terms for workers or workplaces and occupational stress, physical activity, sedentary behaviour, diet, nutrition and lifestyle intervention. Supplementary searches used narrower exposure-specific combinations, together with reference-list screening, citation chasing, publisher records and related-article searching. The complete search strategy is provided in Supplementary Methods S1.

Screening was conducted at title/abstract and full-text stages. Multiple reports from the same cohort or intervention were linked so that overlapping populations were not treated as independent evidence. Because the original record-level screening file was unavailable, final PRISMA identification and exclusion counts will be inserted after completion of the reconstructed screening log.

\subsection{Data extraction and critical appraisal}
A structured extraction framework captured occupational population, design, sample size, exposure or intervention, MetS outcome, principal effect estimate and methodological limitations. Adjusted estimates were preferred for observational studies, and between-group effects were prioritised for interventions where available.

Methodological appraisal used design-specific Joanna Briggs Institute tools for randomised trials, cohort studies, analytical cross-sectional studies and quasi-experimental studies. Appraisal findings were used to qualify interpretation rather than converted into a numerical quality score. Item-level results are reported in Supplementary Tables S2--S5.

\subsection{Synthesis}
Substantial heterogeneity in populations, exposures, interventions, MetS definitions and effect measures precluded a meaningful overall meta-analysis. Evidence was therefore synthesised narratively by domain, with greater interpretive weight given to randomised and prospective evidence than to cross-sectional or uncontrolled findings.

Formal publication-bias analysis was not undertaken because no pooled meta-analysis was performed. A formal GRADE assessment was also not applied; confidence was judged in relation to study design, directness, consistency and JBI appraisal findings.

\section{Results}

\subsection{Evidence base}
The current synthesis comprised 34 primary studies published between 2015 and 2026. Six were randomised trials, six were prospective or longitudinal cohorts, 17 were analytical cross-sectional studies and five were quasi- or non-randomised interventions. The evidence spanned office and white-collar workers, police and security personnel, drivers, manufacturing and petrochemical workers, healthcare staff, university employees, military personnel, heavy-industry workers and participants in workplace health-promotion programmes.

\section*{Table 1. Overview of the evidence base}

\begin{tabular}{p{0.24\textwidth} p{0.28\textwidth} p{0.40\textwidth}}
\toprule
Evidence domain & Main study designs & Overall pattern \\
\midrule
Occupational stress & Prospective cohort; analytical cross-sectional & Prospective studies linked high or increasing work stress with later MetS risk; cross-sectional findings were more mixed and often component-specific. \\
Diet and workplace nutrition & Analytical cross-sectional; randomised and quasi-experimental intervention & Healthier dietary patterns generally aligned with more favourable metabolic profiles, while workplace dietary interventions improved several MetS components. \\
Physical activity and sedentary behaviour & Cohort; analytical cross-sectional; randomised intervention & Leisure-time physical activity was generally associated with lower MetS risk, whereas prolonged sedentary exposure was associated with higher risk. \\
Multicomponent workplace interventions & Randomised and quasi-experimental intervention & Programmes combining diet, activity, counselling, self-monitoring or digital support generally improved anthropometric or metabolic outcomes. \\
\bottomrule
\end{tabular}

\vspace{0.5em}
\noindent\textit{Note:} Several studies contributed evidence to more than one domain. Detailed characteristics for all 34 studies are provided in Supplementary Table S1.
\section*{Table 2. Selected source-verified findings}

\begin{tabular}{p{0.19\textwidth} p{0.27\textwidth} p{0.45\textwidth}}
\toprule
Study & Exposure/intervention & Principal finding \\
\midrule
Garbarino \& Magnavita (2015) & Work stress & High work stress was associated with incident MetS (adjusted OR 2.68, 95\% CI 1.08--6.70). \\
Yamaguchi et al. (2018) & Change in job demands & Low-to-high job demands were associated with incident MetS (adjusted OR 3.27, 95\% CI 1.46--7.34). \\
Kuwahara et al. (2016) & Leisure-time exercise & Compared with the lowest category, mixed moderate/vigorous exercise was associated with lower MetS incidence (HR 0.74, 95\% CI 0.62--0.89 for 7.5--<16.5 MET-h/week). \\
Browne et al. (2017) & Physical activity & Active workers had lower odds of MetS than sedentary workers (OR 0.52, 95\% CI 0.27--0.98). \\
Motuma et al. (2020) & Sedentary behaviour and diet & Sedentary time $\geq$8 h/day was associated with higher MetS prevalence (APR 2.29, 95\% CI 1.88--2.80); high fruit/vegetable intake showed an inverse association. \\
Haufe et al. (2019) & Telemonitored exercise & Exercise reduced MetS severity versus control (between-group difference $-0.26$, 95\% CI $-0.35$ to $-0.16$). \\
Ma et al. (2025) & Canteen-based dietary intervention & The intervention improved fasting glucose, total and LDL cholesterol, waist circumference, BMI and HDL cholesterol. \\
Bogale et al. (2025) & Lifestyle education & Waist circumference decreased by 5.33 cm and systolic blood pressure by 6.96 mmHg, with additional lipid improvements. \\
Ayeltigah et al. (2026) & Work stress in military personnel & Work stress was associated with higher odds of MetS (OR 6.472, 95\% CI 1.535--27.292), although the estimate was imprecise. \\
Ryu et al. (2017) & Targeted workplace health promotion & Waist circumference and fasting glucose improved, while MetS prevalence declined from 41.5\% to 31.7\%. \\
\bottomrule
\end{tabular}

\vspace{0.5em}
\noindent\textit{Abbreviations:} APR, adjusted prevalence ratio; BMI, body mass index; CI, confidence interval; HR, hazard ratio; MetS, metabolic syndrome; OR, odds ratio.

\section*{Table 3. Methodological appraisal by design}

\begin{tabular}{p{0.24\textwidth} p{0.08\textwidth} p{0.24\textwidth} p{0.36\textwidth}}
\toprule
Design & n & Main strengths & Main limitations \\
\midrule
Randomised controlled trials & 6 & Stronger causal inference; standardised outcome assessment & Allocation concealment, masking, attrition and analysis by randomised group were incompletely reported in several trials. \\
Prospective/cohort studies & 6 & Temporal ordering; adjustment for major confounders; standardised MetS assessment & Incomplete reporting of attrition or missing-data handling in several studies. \\
Analytical cross-sectional studies & 17 & Broad occupational coverage; generally standardised MetS criteria & No temporal ordering; reverse causation remains plausible. \\
Quasi-/non-randomised interventions & 5 & Real-world workplace implementation & Non-randomised allocation and, in several studies, absence of a concurrent control group. \\
\bottomrule
\end{tabular}

\vspace{0.5em}
\noindent\textit{Note:} Design-specific Joanna Briggs Institute appraisal was used to qualify confidence in the narrative synthesis; no summed quality score was applied.

\subsection{Occupational stress}
The clearest temporal evidence came from two prospective cohorts. Among police officers, high work stress was associated with incident MetS over five years (adjusted OR 2.68, 95\% CI 1.08--6.70) \citep{Garbarino2015}. In Japanese manufacturing workers, moving from low to high job demands was associated with a greater likelihood of developing MetS over three years (adjusted OR 3.27, 95\% CI 1.46--7.34) \citep{Yamaguchi2018}.

Cross-sectional findings were less uniform. Security-guard data linked occupational stressors with adverse glucose, lipid and blood-pressure profiles \citep{Jovanovic2020}, and work stress was associated with MetS among Ghanaian military personnel, although the estimate was imprecise \citep{Ayeltigah2026Military}. By contrast, Zhang et al. found no overall adjusted association between occupational stress and MetS in petrochemical workers; significant findings were limited to selected stress dimensions and individual MetS components \citep{Zhang2024}. Taken together, the prospective evidence is more persuasive than the cross-sectional literature, which appears sensitive to how stress and MetS are measured.

\subsection{Diet and workplace nutrition}
Dietary findings were broadly consistent in direction but varied in strength. In office workers, low fruit intake was associated with MetS \citep{Alavi2015}. Among Iranian health workers, poorer diet moderation was associated with MetS and several components \citep{Sharifi2025}. In a large Madrid workforce, greater adherence to a Mediterranean dietary pattern was associated with lower odds of MetS \citep{ReinosoBarbero2026}.

Intervention studies suggest that the workplace food environment is modifiable. A randomised canteen-based programme improved fasting glucose, total and LDL cholesterol, HDL cholesterol, waist circumference and body mass index \citep{Ma2025}. A Korean worksite dietary programme also improved multiple MetS indicators among higher-risk workers \citep{Kim2016Dietary}. These findings support dietary modification as a practical workplace strategy, although observational dietary studies remain vulnerable to residual confounding and self-report error.

\subsection{Physical activity and sedentary behaviour}
Physical activity showed the most consistent pattern across study designs. In a prospective cohort of more than 22,000 workers, moderate-to-vigorous leisure-time exercise was associated with lower MetS incidence \citep{Kuwahara2016}. Among workers in sedentary occupations, meeting physical-activity recommendations was associated with lower odds of MetS (OR 0.52, 95\% CI 0.27--0.98) and several individual components \citep{Browne2017}. Taiwanese workers with high leisure-time activity similarly had lower odds of MetS, while commuting activity showed no clear association \citep{Huang2017Taiwan}.

Sedentary exposure showed the opposite pattern. Among working adults in Ethiopia, sitting for at least eight hours per day was associated with more than twice the prevalence of MetS (APR 2.29, 95\% CI 1.88--2.80) \citep{Motuma2020}. Intervention evidence reinforced the potential importance of activity: telemonitoring-supported exercise reduced MetS severity \citep{Haufe2019}, while T'ai Chi and structured workplace activity programmes improved selected metabolic outcomes \citep{Choi2017,Fang2019,ParkHwang2024}.

\subsection{Multicomponent workplace interventions}
Interventions that combined diet, exercise, counselling, self-monitoring or digital support generally moved outcomes in a favourable direction. A voluntary worksite weight-loss programme reduced MetS prevalence in both women and men \citep{Earnest2015}. A targeted workplace health-promotion programme reduced waist circumference and fasting glucose and lowered MetS prevalence from 41.5\% to 31.7\% \citep{Ryu2017}. An internet-based lifestyle programme improved selected cardiometabolic risk factors \citep{Kim2015Internet}, and a web-based high-intensity interval-training programme with calorie restriction produced sustained improvements in weight, blood pressure and triglycerides \citep{So2025}.

The strongest recent intervention evidence came from randomised programmes. In Ethiopian bank employees, a cluster-randomised lifestyle-education intervention reduced waist circumference by 5.33 cm and systolic blood pressure by 6.96 mmHg, alongside favourable lipid changes \citep{Bogale2025}. Across interventions, however, heterogeneity in programme content, intensity and follow-up limits direct comparison.

\section{Discussion}

This review suggests that workplace metabolic risk is best understood as the product of both occupational exposures and modifiable health behaviours rather than as a consequence of any single factor. Across the evidence base, three findings were especially consistent: adverse work stress carried the strongest temporal signal when examined prospectively; greater leisure-time physical activity and lower sedentary exposure were associated with more favourable metabolic profiles; and workplace lifestyle programmes could improve several MetS components.

The work-stress evidence deserves careful interpretation. The two prospective studies provide meaningful temporal support because exposure preceded incident MetS \citep{Garbarino2015,Yamaguchi2018}. Yet the cross-sectional literature was mixed, and one large petrochemical study found no overall adjusted association with MetS \citep{Zhang2024}. This inconsistency argues against treating work stress as a single uniform exposure. Job demands, effort--reward imbalance, control, social support, shift structure and recovery opportunities may influence metabolic risk through different pathways, and their effects may depend on concurrent behaviours such as sleep, diet and physical activity.

The physical-activity evidence was more coherent. Prospective, cross-sectional and intervention studies all pointed towards a beneficial role for leisure-time exercise, while prolonged sedentary exposure was associated with greater MetS risk \citep{Kuwahara2016,Browne2017,Huang2017Taiwan,Motuma2020}. Importantly, occupational activity and commuting activity did not consistently show the same relationship as leisure-time exercise. This distinction matters for workplace policy: simply having a physically demanding job is not equivalent to structured or discretionary physical activity that improves cardiorespiratory fitness.

Dietary evidence also supports intervention at the workplace level, but the observational findings are less secure because diet is difficult to measure and is socially patterned. The association between Mediterranean-style eating and lower MetS prevalence is biologically plausible, while the randomised canteen intervention provides more direct evidence that changing the food environment can improve metabolic markers \citep{ReinosoBarbero2026,Ma2025}. For occupational health practice, this suggests that individual counselling may be more effective when combined with environmental changes that make healthier choices easier during the working day.

The intervention literature is encouraging but heterogeneous. Randomised programmes demonstrated meaningful improvements in MetS severity, waist circumference, blood pressure, glucose regulation and lipids \citep{Haufe2019,Ma2025,Bogale2025}. Quasi-experimental studies often showed similar directions of effect, but their lack of randomisation or concurrent controls reduces causal certainty. The practical lesson is therefore not that one workplace programme is universally effective, but that metabolic risk appears modifiable when interventions are sufficiently intensive, sustained and tailored to the workplace context.

This review has several strengths. It focuses specifically on employed populations, integrates psychosocial and behavioural determinants within one framework, and distinguishes the evidential weight of randomised, prospective and cross-sectional studies. Study-level effect estimates were source-checked, and design-specific JBI appraisal was used to qualify interpretation. Detailed extraction and appraisal results are retained in the supplementary material rather than crowding the main text.

The limitations are equally important. The original screening spreadsheet from the undergraduate review was unavailable, which required a fresh search and re-screening process. Final PRISMA identification and exclusion counts remain dependent on completion of the reconstructed record-level screening log. The review was not prospectively registered, no formal publication-bias analysis was performed, and no GRADE certainty framework was applied. In addition, substantial heterogeneity in occupation, exposure definition, MetS criteria, intervention design and effect measure precluded an overall meta-analysis.

For practice, the evidence favours a multidimensional approach. Occupational-health programmes that address psychosocial stress while also improving food environments, increasing opportunities for physical activity, reducing prolonged sitting and identifying workers at elevated cardiometabolic risk may be more useful than single-component interventions. For research, stronger prospective studies and pragmatic workplace trials are needed, particularly in low- and middle-income settings and occupations that remain under-represented.

\section{Conclusion}

Metabolic syndrome among working adults is associated with a combination of occupational and behavioural factors. Prospective evidence supports adverse work stress as a relevant risk marker, while physical activity, sedentary behaviour and diet show broadly consistent associations with metabolic health. Workplace-based lifestyle interventions can improve several MetS components, although the strength of inference varies by design. The workplace is therefore both a setting in which metabolic risk accumulates and a practical setting in which it can be modified.

\section*{Registration and protocol}
This review update was not prospectively registered. No prospectively registered protocol was available; the revised eligibility criteria, search framework, screening rules, extraction structure and appraisal workflow are archived in the project repository.

\section*{Funding}
No external funding was received for this systematic review.

\section*{Competing interests}
None declared.

\section*{Ethics approval}
Ethics approval was not required because this systematic review synthesises previously published studies and did not involve collection of new participant-level data.

\section*{Data availability}
Study-level extraction tables, appraisal files, search documentation and reproducibility materials are available at \url{https://github.com/OG-Kyere/workplace-stress-metabolic-syndrome-systematic-review}.

\section*{Author contributions}
Gideon Ofosu Kyere conceived the review update, developed the eligibility framework, conducted the updated search and screening, extracted and synthesised the evidence, performed methodological appraisal, and drafted and revised the manuscript.

\section*{Use of artificial intelligence}
Generative artificial-intelligence tools were used for editorial and organisational assistance during manuscript development. All eligibility decisions, methodological judgments, extracted study results, interpretations and final manuscript content were reviewed by the author against the underlying sources.

\bibliographystyle{unsrtnat}
\bibliography{references_verified}

\end{document}